\scanpage{4}
\mspage{2}
\[
\left\{\begin{aligned}
\underline{\unsure{f|p|\rho}}_{m,n} &: \Phi_m\ \unsure{\text{---}|\longrightarrow}\ \Phi_n\subset\Phi_m\\
\underline{\rho}_{n,m}(\{\unsure{\alpha_{n'}|\alpha_{n}^{i}}\})&=
\begin{cases}
\alpha_i=i\in\operatorname{Ob}\Phi_n&\unsure{\text{si}\ \ldots}\quad -n\leq\unsure{i|i'}\leq n\\
\alpha_i=n&\text{si}\quad n\leq i\leq m\\
\alpha_i=-n&\text{si}\quad -m\leq i\leq-n
\end{cases}
\end{aligned}\right.
\]
On note le système projectif d'ens ordonnés ainsi obtenu par $\mathbb{\Phi}=(\Phi_n)_{n\in\mathbb N}$
\[
\mathbb{\Phi}=(\Phi_n)_{n\in\mathbb N}\quad\in\operatorname{Pro}(\mathrm{Cat})
\]
\deleted{\unsure{Pour}} Soit $X$ une catégorie, alors pour $n\in\mathbb N$
\[
\underline{\deleted{\Phi}\mathrm{Ch}}_n(X)\overset{\mathrm{déf}}{=}\underline{\mathrm{Hom}}(\Phi_n,X)
\]
est la catégorie des diagrammes de $X$ de la forme
\diagram{\begin{tikzcd}[ampersand replacement=\&,column sep=.52em,scale=.82,transform shape]
(*)\quad A_{-n}\ \cdots \& A_{-2}\arrow[r,dash,"u_{-2}","{\unsure{\text{---}|\to}}"'] \& A_{-1} \& A_0\arrow[l,"u_{-1}"']\arrow[r,"u_0"] \& A_1\arrow[r,draw=none,"{\deleted{\to}}" description] \& A_2\arrow[l,"u_1"']\arrow[r,"u_2"] \& A_3 \& \cdots\arrow[l] \& A_n
\end{tikzcd}}
avec $2n$ flèches bout à bout, \inserted{$u_i$ $(-n\leq i\leq n)$} deux flèches consécutives étant \unsure{toujours|disposées} en sens opposés, \deleted{\unsure{soit}}
\diagram{\begin{tikzcd}[ampersand replacement=\&,column sep=1em]
\text{soit} \& \cdot\arrow[r] \& \cdot \& \cdot\arrow[l] \& {,}\quad\text{soit} \& {} \& \cdot\arrow[l]\arrow[r] \& {} \& .
\end{tikzcd}}
Le foncteur canonique déduit de $\rho_{n,m}:\Phi_m\to\Phi_n$
\[
\rho^{m,n}_{\deleted{\unsure{X}}}:\underline{\deleted{\Phi}\mathrm{Ch}}_n(X)\longrightarrow\underline{\deleted{\Phi}\mathrm{Ch}}_m(Y)\qquad n\leq m
\]
un \deleted{\unsure{les}} ``foncteur de dégénérescence'', consiste à \deleted{l'opération de} \inserted{compléter} un \deleted{suite} diagramme $(*)$, par $m-n$ termes \deleted{à chaque bout,} $A_{-n}$ avec des identités à gauche, et par $m-n$ termes avec des identités à droite, donc pour $n$ pair on trouve
\diagram{\begin{tikzcd}[ampersand replacement=\&,column sep=.32em,scale=.70,transform shape,row sep=.28em]
A_{-n} \& \cdots\arrow[r,"\mathrm{id}"] \& A_{-n}\arrow[r,"\mathrm{id}"] \& A_{-n}\arrow[r] \& A_{-(n-1)} \& \cdots \& \cdots \& A_{n-1} \& A_n\arrow[l]\arrow[r,"\mathrm{id}"] \& A_n \& A_n\arrow[l,"\mathrm{id}"']\arrow[r] \& \cdots \& A_n\\
\& \smash{\underbrace{\hspace{6em}}_{m-n\ \text{termes}}} \& \& \& \& \& \& \& \& \& \smash{\underbrace{\hspace{7em}}_{m-n\ \text{termes}}}
\end{tikzcd}}
C'est un foncteur pleinement fidèle (ce qui correspond au fait que $\rho_{n,m}$ est surjectif sur les objets et sur les flèches).

% archive sheet number: 3
